\documentclass[../../main.tex]{subfiles}
\graphicspath{{\subfix{./image/}}{\subfix{../../image/}}}
\begin{document}
% SOURCE p=194; printed=179; object=ch05-notes-title; region=top
\section{注记与参考文献}\label{section:ch05-notes-references}\label{section:ch05-notes}
% SOURCE p=194; printed=179; object=ch05-notes-functional-analysis-books; region=upper
泛函分析是一门庞大的学科，本章仅触及其表面。希望进一步学习的读者，可以从 Reed 与 Simon\cite{folland-bib-112}以及 Rudin\cite{folland-bib-126}入手；也应熟悉 Dunford 与 Schwartz\cite{folland-bib-035}以及 Yosida\cite{folland-bib-163}的专著。
% SOURCE p=194; printed=179; object=ch05-notes-functional-analysis-history; region=upper
泛函分析起源于一些经典问题，特别是微分方程与积分方程理论中的问题。对具体无限维函数空间的系统研究，大约始于1907年前后 F.~Riesz、Fréchet、Schmidt、Helly 等人的工作；抽象赋范向量空间的概念则于1920年前后出现在若干作者的论文中。随后十年的研究在 Banach 的经典著作\cite{folland-bib-009}中达到一个高峰，标志着泛函分析成为一门成熟学科。详细的历史叙述见 Dieudonné\cite{folland-bib-033}以及 Dunford 与 Schwartz\cite{folland-bib-035}中的注记。
% SOURCE p=194; printed=179; object=ch05-bochner-integral-definition; region=upper-middle
\cref{section:ch05-normed-spaces}：
\begin{definition}[Bochner 积分]\label{definition:ch05-bochner-integral}
\cref{exercise:ch05-vector-integration}中建立的向量值函数积分，称为 \textbf{Bochner 积分}。
\end{definition}
% SOURCE p=194; printed=179; object=ch05-bochner-separable-range; region=middle
可以删除 $\mathscr Y$ 可分的假设，但此时必须要求 $L_{\mathscr Y}^1$ 中的函数在某个零集上修改后具有可分值域。更详细的论述见 Cohn\cite{folland-bib-027}或 Yosida\cite{folland-bib-163}。
% SUPPLEMENT: id=v14-ch05-bochner-range-details; parent=ch05-bochner-separable-range; kind=proof+motivation+hypotheses; origin=independent-derivation
\begin{proof}
\textbf{可分值域归约（重构新增）。}设 $\mathscr Y$ 为 Banach 空间，$f$ 有一个 Borel 可测代表，$\int\|f\|\,\mathrm d\mu<\infty$，且存在可测零集 $N$，使 $f(X\setminus N)$ 是可分子集。在 $N$ 上将 $f$ 置零，所得 $\widetilde f$ 仍可测。取该值域的可数稠密子集 $D$，令 $\mathscr Z=\overline{\operatorname{span}(D\cup\{0\})}$。$D$ 中向量的有限有理系数线性组合（复情形取 $\mathbb Q+\mathrm i\mathbb Q$ 系数）构成可数集；逐个逼近任意有限组合的系数，说明此集在 $\mathscr Z$ 中稠密。因此 $\mathscr Z$ 可分。它又是 $\mathscr Y$ 的闭子空间，故完备，并包含 $\widetilde f$ 的全部值域。

相对开集的原像表明 $\widetilde f:X\to\mathscr Z$ 可测。于是可以在 $\mathscr Z$ 中直接应用\cref{exercise:ch05-vector-integration}，得到积分，再经等距包含映射视为 $\mathscr Y$ 的元素。此结果与所选可分闭子空间无关：两种选择 $\mathscr Z_1,\mathscr Z_2$ 都包含于可分 Banach 子空间 $\overline{\mathscr Z_1+\mathscr Z_2}$，习题中的有界算子与积分交换性质作用于两个包含映射，使两种积分在该共同空间中相同。更换几乎处处相等的可测代表也不改变积分，因为差的积分范数不超过差的 $L^1$ 范数，后者为零。
\end{proof}
\begin{remark}
\textbf{动机与逐条件说明（重构新增）。}积分构造实际只使用某一个函数可到达的方向，因此可把“整个值域空间可分”减弱为“该函数本质上具有可分值域”。可数稠密集使习题中的可测简单函数逼近可行；范数可积控制积分误差；Banach 完备性保证近似积分的极限留在值域空间中。这里明确使用可测零集和可测代表，不能在不完备测度下任意修改非可测的零集子集后仍直接声称可测。此归约不要求 $\mu$ 有限或 $\sigma$ 有限。
\end{remark}
% END SUPPLEMENT: id=v14-ch05-bochner-range-details
% SOURCE p=194; printed=179; object=ch05-weak-integral-definition; region=middle
\begin{definition}[弱可积与弱积分]\label{definition:ch05-weak-integral}
向量值积分还有另一种定义方法。设 $(X,\mathscr M,\mu)$ 为测度空间，$\mathscr Y$ 为拓扑向量空间，其连续线性泛函能够分离点。函数 $f:X\to\mathscr Y$ 称为\textbf{弱可积}的，若：
\begin{enumerate}
\item $\phi\circ f\in L^1(\mu)$ 对每个 $\phi\in\mathscr Y^*$ 成立；
\item 存在 $y\in\mathscr Y$（它必定唯一），使
\begin{align*}
\int\phi\circ f\,\mathrm d\mu=\phi(y)\qquad(\phi\in\mathscr Y^*).
\end{align*}
\end{enumerate}
此时定义 $\int f\,\mathrm d\mu=y$。
\end{definition}
% SOURCE p=194; printed=179; object=ch05-weak-bochner-comparison-original; region=middle
若 $\mathscr Y$ 为可分 Banach 空间，这种积分概念与 Bochner 积分一致。见 Yosida\cite{folland-bib-163}以及 Rudin\cite{folland-bib-126}。
% SUPPLEMENT: id=v14-ch05-weak-integral-details; parent=ch05-weak-integral-definition; kind=proof+source-clarification+motivation+hypotheses; origin=independent-derivation
\begin{proof}
\textbf{唯一性与线性补写（重构新增）。}若 $y,z$ 都满足定义中的等式，则每个 $\phi\in\mathscr Y^*$ 都满足 $\phi(y-z)=0$。连续对偶分离点，故 $y=z$。若 $f,g$ 弱可积，积分分别为 $y,z$，则对任意标量 $a,b$ 及任意 $\phi$，标量函数 $\phi\circ(af+bg)=a\phi\circ f+b\phi\circ g$ 可积，且积分为 $\phi(ay+bz)$。因此 $af+bg$ 弱可积，其积分为 $ay+bz$。
\end{proof}
\begin{remark}
\textbf{动机、条件与原句范围（重构新增）。}弱积分用全部连续标量测量来规定向量的积分。第一条件使每一个标量积分有定义；第二条件要求所有这些数确实由 $\mathscr Y$ 中同一个向量表示，不能仅检验第一条件便省去存在性。连续对偶分离点只负责表示向量的唯一性；若它不能分离点，零函数也可能有不唯一的“积分”。这里不要求范数、可分性或完备性，因此定义适用于更广的拓扑向量空间。

原文紧随定义的“一致”应区别两层意思：对 Bochner 可积函数，两种积分取同一值；但可分 Banach 空间上的弱可积函数未必 Bochner 可积，因此若把原句理解为两类可积函数相同，就不成立。下一例给出反例，并补证正确的单向包含。
\end{remark}
% END SUPPLEMENT: id=v14-ch05-weak-integral-details
% SUPPLEMENT: id=v14-ch05-weak-not-bochner-example; parent=ch05-weak-bochner-comparison-original; kind=example+proof+source-correction+motivation+hypotheses; origin=independent-derivation
\begin{example*}[弱可积但非 Bochner 可积（重构新增）]\phantomsection\label{example:ch05-weak-not-bochner}
在 $X=\mathbb N$ 上取计数测度，在可分 Hilbert 空间 $\mathscr Y=\ell^2(\mathbb N)$ 中令 $f(n)=e_n/n$，其中 $e_n$ 为标准单位向量。则 $f$ 弱可积，而不 Bochner 可积。
\end{example*}
\begin{proof}
任意 $\phi\in\mathscr Y^*$ 由\cref{theorem:ch05-riesz-hilbert}表示为 $\phi(x)=\langle x,a\rangle$，其中 $a=(a_n)\in\ell^2$。Schwarz 不等式给
\begin{align*}
\sum_{n=1}^{\infty}|\phi(f(n))|
=\sum_{n=1}^{\infty}\frac{|a_n|}{n}
\leqslant\|a\|_2\left(\sum_{n=1}^{\infty}\frac1{n^2}\right)^{1/2}<\infty.
\end{align*}
令 $v=(1/n)_{n\geqslant1}\in\ell^2$。部分和 $v_N=\sum_{n=1}^Ne_n/n$ 在范数下趋于 $v$，故由 $\phi$ 连续，$\sum_n\phi(f(n))=\phi(v)$。这验证了弱可积的两项要求。另一方面，$\int\|f\|_2\,\mathrm d\mu=\sum_n1/n=\infty$，故 $f$ 不属于 Bochner 可积函数类。此例所有函数都可测，值域本身可数，因而不是可测性或可分性缺失造成的。

现在设 $f$ 是取值于 Banach 空间的 Bochner 可积函数。采用前面的可分子空间归约后，由\subcref{exercise:ch05-vector-integration}{item:ch05-vector-integral-operator}，任意连续线性泛函 $\phi$ 都满足
\begin{align*}
\int|\phi\circ f|\,\mathrm d\mu\leqslant\|\phi\|\int\|f\|\,\mathrm d\mu<\infty,
\qquad
\int\phi\circ f\,\mathrm d\mu=\phi\left(\int f\,\mathrm d\mu\right).
\end{align*}
Banach 空间的连续对偶分离点，故这就是弱积分，且积分值一致。在可分 Banach 空间中，若另有 Borel 可测性和范数可积性，则直接满足\cref{exercise:ch05-vector-integration}的条件，因而 Bochner 可积；不能只用逐泛函的标量可积性代替范数可积性。
\end{proof}
\begin{remark}
\textbf{条件差异（重构新增）。}反例把不可和的长度 $1/n$ 分散到互相正交的方向；每个固定的连续泛函的系数属于 $\ell^2$，所以仍产生绝对可和的标量级数。对“每个测量分别可积”的要求较弱，它没有给出长度函数的积分控制。即使对每个子集分别要求存在弱积分，此例也仍成立：对子集 $E\subseteq\mathbb N$，相应积分向量为 $(\chi_E(n)/n)_n\in\ell^2$，同一计算即可验证。
\end{remark}
% END SUPPLEMENT: id=v14-ch05-weak-not-bochner-example
% SOURCE p=194; printed=179; object=ch05-notes-open-closed-history; region=middle-lower
\cref{section:ch05-baire-category}：开映射定理和闭图像定理归于 Banach\cite{folland-bib-009}。关于开映射定理与 Tietze 延拓定理的证明之间的关系，Grabiner\cite{folland-bib-058}给出了有趣的评论。
% SOURCE p=194; printed=179; object=ch05-notes-uniform-boundedness-history; region=lower
本书表述的一致有界原理归于 Banach 与 Steinhaus\cite{folland-bib-010}；不过，定理的第二部分，即若 $\mathscr X$ 为 Banach 空间且 $\sup_{T\in\mathscr A}\|Tx\|<\infty$ 对每个 $x\in\mathscr X$ 成立，则 $\sup_{T\in\mathscr A}\|T\|<\infty$，此前已经用 Dieudonné\cite{folland-bib-033}所称的“滑动驼峰法”证明。这一颇为精巧而初等的论证近年来在很大程度上受到忽视，但可以在 Hennefeld\cite{folland-bib-071}中找到现代的叙述。
% SUPPLEMENT: id=v14-ch05-gliding-hump-proof; parent=ch05-notes-uniform-boundedness-history; kind=proof+motivation+hypotheses; origin=independent-derivation
\begin{proof}
\textbf{滑动驼峰式补证（重构新增）。}以下独立证明原段所述的一致有界结论，不把此推导归作所引文献的逐字证明。设 $\mathscr A\subseteq L(\mathscr X,\mathscr Y)$，$\mathscr X$ 为 Banach 空间，且每个 $x$ 都满足 $\sup_{T\in\mathscr A}\|Tx\|<\infty$。反设 $\sup_{T\in\mathscr A}\|T\|=\infty$。递归选择 $T_n\in\mathscr A$ 和 $x_n\in\mathscr X$。已选到 $n-1$ 时，令
\begin{align*}
r_n=\frac{2^{-n}}{1+\max_{1\leqslant j<n}\|T_j\|},
\qquad s_{n-1}=\sum_{j<n}x_j,
\end{align*}
其中空集的最大值在此约定为0。选 $\|T_n\|>4n/r_n$。由算子范数定义，存在 $v_n$ 满足 $\|v_n\|\leqslant r_n$ 且 $\|T_nv_n\|>2n$。三角不等式给
\begin{align*}
2\|T_nv_n\|\leqslant
\|T_n(s_{n-1}+v_n)\|+\|T_n(s_{n-1}-v_n)\|.
\end{align*}
故可取 $x_n=v_n$ 或 $x_n=-v_n$，使 $\|T_ns_n\|>2n$，其中 $s_n=s_{n-1}+x_n$。

由于 $\sum_n\|x_n\|\leqslant\sum_n2^{-n}<\infty$，$\mathscr X$ 完备性保证 $s_n\to x\in\mathscr X$。固定 $n$，当 $k>n$ 时有 $r_k\leqslant2^{-k}/(1+\|T_n\|)$，所以 $T_n$ 的连续性和尾和估计给
\begin{align*}
\|T_n(x-s_n)\|\leqslant\|T_n\|\sum_{k>n}r_k
\leqslant\sum_{k>n}2^{-k}=2^{-n}.
\end{align*}
于是 $\|T_nx\|>2n-2^{-n}\to\infty$，这与在同一个 $x$ 处的逐点有界性矛盾。一致有界性得证。
\end{proof}
\begin{remark}
\textbf{动机与逐条件说明（重构新增）。}第 $n$ 步制造一个被 $T_n$ 放大的小扰动，再让后面全部扰动对已选算子的总影响足够小。正负号的选择防止已有部分和抵消新扰动。每个算子有界使后续尾部可控；整族范数无上界使任意小扰动仍能产生大像；定义域完备使这些扰动累加为单个可测试的向量。值域只需赋范，不需完备。这与 Baire 证明走的是不同路径，却证明同一结论。
\end{remark}
% END SUPPLEMENT: id=v14-ch05-gliding-hump-proof
% SOURCE p=194; printed=179; object=ch05-notes-unbounded-maps; region=lower
当 $\mathscr X$ 不完备时，容易构造两个赋范向量空间之间的无界线性映射 $T:\mathscr X\to\mathscr Y$（见\cref{exercise:ch05-weighted-incomplete,exercise:ch05-derivative-incomplete}）；然而，当 $\mathscr X$ 完备时，若不使用选择公理，几乎不可能作出这种构造。标准方法如下。
% SOURCE p=194; printed=179; object=ch05-unbounded-completion-extension; region=lower
\begin{proposition}\label{proposition:ch05-unbounded-completion-extension}
从一个无界线性映射 $T:\mathscr X_0\to\mathscr Y$ 出发，其中 $\mathscr X_0$ 不完备，令 $\mathscr X$ 为 $\mathscr X_0$ 的完备化。则可以把 $T$ 延拓到 $\mathscr X$ 上。
\end{proposition}
% SOURCE p=194; printed=179; object=ch05-unbounded-completion-extension-proof; region=bottom
\begin{proof}
选取 $\mathscr X_0$ 的一个基 $\{u_\alpha\}_{\alpha\in A}$。
\begin{definition*}[本段中的基]\phantomsection\label{definition:ch05-hamel-basis-notes}
这里的基是代数意义的基，即每个 $x\in\mathscr X_0$ 是这些 $u_\alpha$ 的\textbf{有限}线性组合。
\end{definition*}
% SUPPLEMENT: id=v14-ch05-hamel-local-remark; parent=ch05-unbounded-completion-extension-proof; kind=motivation+hypotheses; origin=independent-derivation
\begin{remark}
\textbf{局部定义的条件说明（重构新增）。}本段选代数基，是为了逐个新增方向指定延拓的值。“基”同时要求线性无关与张成；张成保证每个向量有有限展开，线性无关保证该展开唯一。这里只用有限组合，因此不需要对无穷级数作任何连续性承诺。原括号强调有限性，并未取消基的线性无关要求；它与上一节允许范数收敛无限展开的标准正交基不同。
\end{remark}
% END SUPPLEMENT: id=v14-ch05-hamel-local-remark
% SOURCE p=194; printed=179; object=ch05-unbounded-completion-extension-proof-continuation; region=bottom
把它扩充为 $\mathscr X$ 的基 $\{u_\alpha\}_{\alpha\in B}$，其中 $B\supset A$。（这里用到了选择公理。）令 $\mathscr M$ 为 $\{u_\alpha\}_{\alpha\in B\setminus A}$ 的线性张成空间。于是每个 $x\in\mathscr X$ 都可以唯一地写为 $x=x_0+x_1$，其中 $x_1\in\mathscr X_0$ 且 $x_1\in\mathscr M$。令 $T(x_0+x_1)=Tx_0$，便将 $T$ 延拓到了 $\mathscr X$。
\end{proof}
% SUPPLEMENT: id=v14-ch05-unbounded-extension-details; parent=ch05-unbounded-completion-extension; kind=proof-detail+source-correction+motivation+hypotheses; origin=independent-derivation
\begin{remark}
\textbf{原记号修正与定义范围（重构新增）。}原证明分解后的第一个 $x_1\in\mathscr X_0$ 应为 $x_0\in\mathscr X_0$，第二个仍为 $x_1\in\mathscr M$。这里“基”沿用线性代数中“线性无关且张成”的完整含义；原括号强调有限组合，以区别 Hilbert 空间的无穷正交展开，并非把任意生成族都叫基。延拓只要求线性，其限制仍无界，所以绝不是连续延拓。
\end{remark}
\begin{proof}
\textbf{基的扩充及延拓补写（重构新增）。}任意向量空间的线性无关子集按包含关系排序，每条链的并仍线性无关：任何有限线性关系只涉及有限多个向量，它们同属链中某个成员。空链也有空集作上界。由\hyperref[lemma:ch00-zorn]{Zorn 引理}存在极大线性无关子集；若其线性张成不是整个空间，可再加入一个不在张成空间的向量，矛盾。因此它是代数基。先在 $\mathscr X_0$ 中作此构造，再在 $\mathscr X$ 中对包含已选基的线性无关子集作相同构造，得到所需扩充基。

把任意 $x$ 的唯一有限基展开分成 $A$ 内和 $B\setminus A$ 内两部分，得到 $x=x_0+x_1$。两部分张成空间之交为零，因为交中的非零表示将给出整组基的非平凡有限线性关系。因此分解唯一，且分量映射 $P_0x=x_0$ 线性。定义 $\widetilde T=T\circ P_0$，它线性，且 $x\in\mathscr X_0$ 时 $P_0x=x$，所以确实延拓原 $T$。若 $\widetilde T$ 有界，某个 $C$ 满足 $\|Tx\|=\|\widetilde Tx\|\leqslant C\|x\|$ 对所有 $x\in\mathscr X_0$ 成立，与原 $T$ 无界矛盾。故所得延拓仍无界。
\end{proof}
\begin{remark}
\textbf{动机与逐条件说明（重构新增）。}代数基把延拓化为在新增方向上任意指定值，这里全部指定为零。选择原理负责基的存在与扩充；完备化保证新的定义域为 Banach 空间，并等距保留原范数，使原有的无界性不会消失。值域无需完备，补空间 $\mathscr M$ 也不要求闭，因为整个构造不使用拓扑极限。给定无界 $T$ 后，纯代数延拓本身甚至可在任意包含 $\mathscr X_0$ 的向量空间中进行；原文的“不完备”服务于先获得具体的无界例子。这里没有把历史叙述“几乎不可能”误写成某个未经证明的选择公理等价命题。
\end{remark}
% END SUPPLEMENT: id=v14-ch05-unbounded-extension-details
% SOURCE p=195; printed=180; object=ch05-notes-tvs-history; region=top
\cref{section:ch05-topological-vector-spaces}：Trèves\cite{folland-bib-150}对拓扑向量空间的一般理论作了易读的论述，并给出许多具体例子。
% SOURCE p=195; printed=180; object=ch05-notes-alaoglu-history; region=upper
Alaoglu 定理最早由 Alaoglu\cite{folland-bib-003}宣布，详细证明见 Alaoglu\cite{folland-bib-004}；它包含此前若干处理特殊情形的结果。Bourbaki\cite{folland-bib-019}也独立发现了这一定理。
% SOURCE p=195; printed=180; object=ch05-notes-hilbert-history; region=upper-middle
\cref{section:ch05-hilbert-spaces}：Hilbert 本人考虑的空间为 $\ell^2(\mathbb N)$；抽象 Hilbert 空间的概念由 von Neumann\cite{folland-bib-154}在量子力学的数学研究中引入。\cref{theorem:ch05-riesz-hilbert}最早由 F.~Riesz\cite{folland-bib-115}在 $L^2$ 空间背景下得到。这是若干以他的名字命名的、关于不同空间上线性泛函的表示定理之一，另外几个为\cref{theorem:ch06-lp-duality,theorem:ch07-riesz-positive,theorem:ch07-riesz-complex}。为避免混淆，本书把“Riesz 表示定理”这一名称保留给后两个关系密切的定理。
% SOURCE p=195; printed=180; object=ch05-notes-physics-inner-product; region=middle
在量子物理文献中，标量积通常写为 $\langle x\mid y\rangle$，并约定它对第二个变量线性、对第一个变量共轭线性。
% SUPPLEMENT: id=v14-ch05-physics-convention-details; parent=ch05-notes-physics-inner-product; kind=motivation+hypotheses; origin=independent-derivation
\begin{remark}
\textbf{约定转换（重构新增）。}若沿用本书第一变量线性的内积，可定义物理约定的 $\langle x\mid y\rangle=\langle y,x\rangle=\overline{\langle x,y\rangle}$。于是第二变量的线性直接来自右端第一变量的线性，第一变量的反线性也相应得到；平方范数、正交关系均不变。需要转换的是含混合内积的共轭位置，不能在同一个证明中只换括号而保留不相容的线性约定。前文补充说明提及 Hilbert 空间表示定理的通用名称；本书原有的定理命名仍按本段所述约定保留。
\end{remark}
% END SUPPLEMENT: id=v14-ch05-physics-convention-details
\end{document}
